\documentclass[11pt]{article}
\usepackage[right=1in,left=1in,top=1.1in,bottom=1.1in]{geometry}
\usepackage{amsmath,amsthm,amssymb}
\usepackage{hyperref}
\hypersetup{colorlinks, citecolor=blue, filecolor=blue, linkcolor=blue, urlcolor=blue}
\usepackage{setspace}
\onehalfspacing
\usepackage{sectsty}
\sectionfont{\large}
\subsectionfont{\normalsize}

\newtheorem{proposition}{Proposition}
\newtheorem{remark}{Remark}

\title{IQ-Learn as Penalized Conditional Log-Likelihood:\\A Tabular Reinterpretation for the DDC Setting}
\author{Pranjal Rawat}
\date{March 2026}

\begin{document}
\maketitle

\section{Preliminaries}

The original IQ-Learn paper (Garg et al., 2021) is designed for model-free imitation learning in continuous environments where the learner has no explicit knowledge of the transition dynamics. In this note we work in the finite discrete setting of DDC models, where transition probabilities $p(s'|s,a)$ are known or consistently estimated from data. This allows us to compute continuation values exactly and rewrite the IQ-Learn chi-squared objective in a form familiar to structural econometricians.

The main finding is that the chi-squared IQ-Learn objective, when evaluated with known transitions, is algebraically equivalent to a penalized conditional log-likelihood. The first term is the standard multinomial logit likelihood that fits expert actions through a softmax policy. The second term penalizes the squared magnitude of the Bellman-implied reward at expert-visited states. The $Q$-function is the single free parameter that determines both the policy (via softmax) and the implied reward (via the inverse Bellman operator). No inner Bellman loop is needed because the smoothed value function is an explicit closed-form function of $Q$. The derivation requires only two algebraic substitutions and is verified numerically on a small MDP with comparisons to NFXP.

Consider a finite-state, finite-action Markov decision process with state space $S$, action set $A(s)$ for each state $s$, transition probabilities $p(s'|s,a)$, per-period reward $r(s,a)$, discount factor $\beta \in (0,1)$, and extreme value scale parameter $\sigma > 0$. Under the additive separability, conditional independence, and extreme value (AS/CI/EV) assumptions, the expected value function (the ``smoothed max function'') is
\begin{equation}
V(s) = \sigma \log\!\left(\sum_{a \in A(s)} \exp\{Q(s,a)/\sigma\}\right),
\label{eq:smoothV}
\end{equation}
where the choice-specific value function $Q$ is the unique fixed point of the smoothed Bellman operator $\Lambda_\sigma$,
\begin{equation}
\Lambda_\sigma(Q)(s,a) = r(s,a) + \beta \sum_{s'} V(s')\, p(s'|s,a).
\label{eq:smoothQ}
\end{equation}
The right-hand side of \eqref{eq:smoothQ} depends on the entire $Q$-table through $V(s')$ in \eqref{eq:smoothV}, so $\Lambda_\sigma$ maps $Q$-functions to $Q$-functions.
The conditional choice probability takes the multinomial logit form
\begin{equation}
\pi(a|s) = \frac{\exp\{Q(s,a)/\sigma\}}{\sum_{a' \in A(s)} \exp\{Q(s,a')/\sigma\}}.
\label{eq:ccp}
\end{equation}

\section{Two Identities}

\begin{proposition}[Advantage equals scaled log choice probability]
\label{prop:identity1}
For any function $Q(s,a)$, define $V(s)$ by \eqref{eq:smoothV} and $\pi(a|s)$ by \eqref{eq:ccp}. Then
\begin{equation}
Q(s,a) - V(s) = \sigma \log \pi(a|s).
\label{eq:identity1}
\end{equation}
\end{proposition}

\begin{proof}
Take the logarithm of \eqref{eq:ccp}.
\[
\log \pi(a|s) = \frac{Q(s,a)}{\sigma} - \log\!\left(\sum_{a'} \exp\{Q(s,a')/\sigma\}\right).
\]
Multiply both sides by $\sigma$.
\[
\sigma \log \pi(a|s) = Q(s,a) - \sigma \log\!\left(\sum_{a'} \exp\{Q(s,a')/\sigma\}\right) = Q(s,a) - V(s). \qedhere
\]
\end{proof}

\noindent In the reinforcement learning literature the quantity $Q(s,a) - V(s)$ is called the advantage function $\mathcal{A}(s,a)$. Identity \eqref{eq:identity1} says the advantage equals the scaled log choice probability under the logit model.

\begin{proposition}[Temporal difference equals implied reward]
\label{prop:identity2}
Define the temporal difference
\begin{equation}
\mathrm{td}(s,a) \;\equiv\; Q(s,a) - \beta \sum_{s'} V(s')\, p(s'|s,a).
\label{eq:td}
\end{equation}
If $Q$ satisfies the Bellman fixed point $Q = \Lambda_\sigma(Q)$, then $\mathrm{td}(s,a) = r(s,a)$.
\end{proposition}

\begin{proof}
At the fixed point, $Q(s,a) = r(s,a) + \beta \sum_{s'} V(s')\, p(s'|s,a)$ by \eqref{eq:smoothQ}. Substituting into \eqref{eq:td} gives $\mathrm{td}(s,a) = r(s,a)$.
\end{proof}

\noindent For a general $Q$ that does not satisfy the Bellman equation, $\mathrm{td}(s,a)$ is the reward that would rationalize $Q$ as a fixed point. We call $\mathrm{td}(s,a)$ the implied reward. The equality $\mathrm{td} = r$ holds only at the Bellman fixed point. But the definition of $\mathrm{td}$ and the identity in Proposition~\ref{prop:identity1} hold for any $Q$-function, not only for $Q^*$. The penalized MLE equivalence derived in the next two sections follows entirely from these definitions and is therefore an algebraic identity valid at every point in the optimization, not an approximation that holds only at convergence.

\section{The IQ-Learn Chi-Squared Objective}

IQ-Learn parameterizes $Q$ directly and optimizes over it. Suppose we observe expert state-action pairs $(s_i, a_i)$ for $i = 1, \ldots, N$. The chi-squared objective (Equation 12 of Garg et al., 2021) is
\begin{equation}
\mathcal{L}(Q) \;=\; -\frac{1}{N}\sum_{i=1}^{N}\bigl[Q(s_i,a_i) - V(s_i)\bigr] \;+\; \frac{1}{4\alpha}\,\frac{1}{N}\sum_{i=1}^{N}\bigl[\mathrm{td}(s_i,a_i)\bigr]^2,
\label{eq:iqlearn}
\end{equation}
where $\alpha > 0$ is a regularization parameter and $V(s)$ is the smoothed value function \eqref{eq:smoothV} evaluated at the current $Q$. The estimator minimizes $\mathcal{L}(Q)$ over $Q$. In the original IQ-Learn paper, the continuation value $\sum_{s'} V(s')\,p(s'|s,a)$ is estimated from sampled next states because the transition probabilities are unknown. In the DDC setting we compute this sum exactly from the known transition matrix.

\section{Derivation}

We now show that minimizing \eqref{eq:iqlearn} is equivalent to maximizing a penalized conditional log-likelihood. The argument proceeds in three steps.

\paragraph{Step 1. Rewrite the first term as scaled negative conditional log-likelihood.}

By Proposition~\ref{prop:identity1}, $Q(s_i,a_i) - V(s_i) = \sigma \log \pi(a_i|s_i)$. Substituting into the first term of \eqref{eq:iqlearn} gives
\begin{equation}
-\frac{1}{N}\sum_{i=1}^{N}\bigl[Q(s_i,a_i) - V(s_i)\bigr] \;=\; -\frac{\sigma}{N}\sum_{i=1}^{N} \log \pi(a_i|s_i).
\label{eq:step1}
\end{equation}
The sum $\sum_{i=1}^{N} \log \pi(a_i|s_i)$ is the conditional log-likelihood of the expert actions given the observed states, under the multinomial logit with $Q$-values as choice-specific utilities. Denote it $\mathrm{LL}(Q)$. The first term equals $-(\sigma/N)\,\mathrm{LL}(Q)$.

\paragraph{Step 2. Rewrite the second term as a squared reward penalty.}

By Proposition~\ref{prop:identity2}, $\mathrm{td}(s_i,a_i)$ is the implied reward at observation $i$. Write $\hat{r}_i \equiv \mathrm{td}(s_i, a_i)$. The second term of \eqref{eq:iqlearn} becomes
\begin{equation}
\frac{1}{4\alpha}\,\frac{1}{N}\sum_{i=1}^{N}\bigl[\mathrm{td}(s_i,a_i)\bigr]^2 \;=\; \frac{1}{4\alpha N}\sum_{i=1}^{N} \hat{r}_i^{\,2}.
\label{eq:step2}
\end{equation}
This is the average squared implied reward at expert-visited state-action pairs.

\paragraph{Step 3. Combine and simplify.}

Substituting \eqref{eq:step1} and \eqref{eq:step2} into \eqref{eq:iqlearn}, the objective is
\[
\mathcal{L}(Q) = -\frac{\sigma}{N}\,\mathrm{LL}(Q) + \frac{1}{4\alpha N}\sum_{i=1}^{N} \hat{r}_i^{\,2}.
\]
Multiplying by $-N/\sigma$ converts minimization of $\mathcal{L}(Q)$ into maximization of
\begin{equation}
\boxed{\;\mathrm{LL}(Q) \;-\; \frac{1}{4\alpha\sigma}\sum_{i=1}^{N} \hat{r}_i^{\,2}\;}
\label{eq:penalizedMLE}
\end{equation}
where $\hat{r}_i = Q(s_i,a_i) - \beta \sum_{s'} V(s')\,p(s'|s_i,a_i)$ is the implied reward at the $i$-th expert observation and $V(s)$ is defined by \eqref{eq:smoothV} as a function of $Q$.

\medskip

Expression \eqref{eq:penalizedMLE} is a penalized conditional log-likelihood. The first term is the multinomial logit conditional log-likelihood with $Q$-values as choice-specific utilities. The second term is an $\ell^2$ penalty on the Bellman-implied reward, weighted by the expert occupancy (since the sum runs over expert observations only). The penalty strength is $1/(4\alpha\sigma)$. The analogy to ridge regression is instructive but inexact. In ridge regression, the penalty is a quadratic function of the parameters. Here, the implied reward $\hat{r}_i = Q(s_i,a_i) - \beta\sum_{s'} V(s')\,p(s'|s_i,a_i)$ is a nonlinear function of $Q$ because $V(s')$ depends on the entire $Q$-table through the logsumexp in \eqref{eq:smoothV}. The penalty is quadratic in $\hat{r}$ but not in $Q$.

\section{Numerical Illustration}
\label{sec:numerical}

We verify the propositions and the penalized MLE equivalence on a small MDP with 2 states and 2 actions. Set $\beta = 0.9$ and $\sigma = 0.5$. The rewards are $r(0,0) = 1$, $r(0,1) = 0.5$, $r(1,0) = 0.2$, $r(1,1) = 0.8$. The transition matrix $p(s'|s,a)$ has four rows:
$p(\cdot|0,0) = (0.7, 0.3)$, $p(\cdot|0,1) = (0.4, 0.6)$, $p(\cdot|1,0) = (0.5, 0.5)$, $p(\cdot|1,1) = (0.3, 0.7)$.

The smoothed Bellman equation $Q = \Lambda_\sigma(Q)$ is solved by value iteration to convergence. Table~\ref{tab:mdp} reports the solution. With $\sigma = 0.5$ the policy is more concentrated than it would be at $\sigma = 1$, assigning 0.762 and 0.749 to the preferred actions.

\begin{table}[ht]
\centering
\caption{MDP solution and verification of Proposition~\ref{prop:identity1}.}
\label{tab:mdp}
\begin{tabular}{ccrrrrr}
\hline
$s$ & $a$ & $Q^*(s,a)$ & $V^*(s)$ & $\pi(a|s)$ & $Q{-}V$ & $\sigma\log\pi$ \\
\hline
0 & 0 & 10.417 & 10.554 & 0.762 & $-0.136$ & $-0.136$ \\
0 & 1 & ~9.837 & 10.554 & 0.238 & $-0.717$ & $-0.717$ \\
1 & 0 & ~9.564 & 10.254 & 0.251 & $-0.691$ & $-0.691$ \\
1 & 1 & 10.110 & 10.254 & 0.749 & $-0.145$ & $-0.145$ \\
\hline
\end{tabular}
\end{table}

\noindent\textbf{Proposition 1.} The last two columns of Table~\ref{tab:mdp} confirm that $Q(s,a) - V(s)$ and $\sigma\log\pi(a|s)$ agree to machine precision. With $\sigma = 0.5$ the advantage $Q - V$ is half the log choice probability rather than equal to it, verifying that the scale parameter enters the identity as stated.

\noindent\textbf{Proposition 2.} Table~\ref{tab:td} reports the temporal difference $\mathrm{td}(s,a) = Q(s,a) - \beta\sum_{s'} V(s')\,p(s'|s,a)$ at the Bellman fixed point. The implied reward equals the true reward to twelve decimal places.

\begin{table}[ht]
\centering
\caption{Temporal difference equals reward at the Bellman fixed point.}
\label{tab:td}
\begin{tabular}{ccrrc}
\hline
$s$ & $a$ & $\mathrm{td}(s,a)$ & $r(s,a)$ & $|\mathrm{td} - r|$ \\
\hline
0 & 0 & 1.0000 & 1.0000 & ${<}10^{-12}$ \\
0 & 1 & 0.5000 & 0.5000 & ${<}10^{-12}$ \\
1 & 0 & 0.2000 & 0.2000 & ${<}10^{-12}$ \\
1 & 1 & 0.8000 & 0.8000 & ${<}10^{-12}$ \\
\hline
\end{tabular}
\end{table}

\noindent\textbf{Penalized MLE equivalence.} Using all four state-action pairs as observations and setting $\alpha = 1$, we compute the IQ-Learn objective \eqref{eq:iqlearn} directly and via the penalized MLE decomposition \eqref{eq:penalizedMLE}. The log-likelihood is $\mathrm{LL}(Q^*) = -3.377$. The sum of squared implied rewards is $\sum \hat{r}_i^2 = 1.930$. The penalty coefficient is $1/(4\alpha\sigma) = 0.5$, so the penalty term is $0.5 \times 1.930 = 0.965$. The IQ-Learn objective evaluates to $\mathcal{L}(Q^*) = 0.543$, and the rescaled quantity $-N/\sigma \cdot \mathcal{L}(Q^*)$ equals $\mathrm{LL}(Q^*) - 0.965 = -4.342$ to machine precision.

\noindent\textbf{Reward shaping invariance.} Take the potential function $h = (2, -1)$. The shaped rewards $r'(s,a) = r(s,a) + h(s) - \beta\sum_{s'}h(s')\,p(s'|s,a)$ range from $-1.25$ to $2.32$, far from the originals. Yet the policy $\pi'(a|s)$ under the shaped reward is identical to $\pi(a|s)$ to machine precision, and $Q'(s,a) = Q(s,a) + h(s)$ holds for all state-action pairs. The log-likelihood $\mathrm{LL}(Q')$ equals $\mathrm{LL}(Q)$, but the implied rewards change. The sum of squared implied rewards rises from $\sum\hat{r}_i^2 = 1.93$ under the original reward to $\sum\hat{r}_i^{\prime 2} = 11.00$ under the shaped reward. The regularizer in \eqref{eq:penalizedMLE} breaks the equivalence by preferring the $Q$-function with the smaller implied reward.

\noindent\textbf{Comparison with NFXP.} We simulate 1000 expert state-action pairs from the true policy and estimate by both NFXP and IQ-Learn. NFXP maximizes the conditional log-likelihood $\mathrm{LL}(\theta)$ subject to the Bellman constraint. With the standard normalization $r(s,0) = 0$ for identification, NFXP recovers the reward differences $r(0,1) - r(0,0) = -0.80$ and $r(1,1) - r(1,0) = 0.36$ (true values $-0.50$ and $0.60$), achieving $\mathrm{LL} = -563.4$. IQ-Learn with $\alpha = 1$ recovers similar reward differences ($-0.51$ and $0.38$) but achieves a lower conditional log-likelihood of $-567.7$. The gap reflects the cost of regularization. NFXP maximizes fit alone, while IQ-Learn trades some fit for a smaller implied reward, achieving $\sum\hat{r}_i^2 = 0.25$ compared to $0.76$ for NFXP. Both estimators recover policies close to the truth.

\section{Discussion}

\subsection{Comparison with NFXP}

In NFXP, the Bellman equation $Q = \Lambda_\sigma(Q)$ is imposed as a hard constraint. The econometrician parameterizes the reward function $r_\theta(s,a) = \sum_k \theta_k \phi_k(s,a)$ and solves the fixed point $Q(\theta) = \Lambda_\sigma(Q(\theta))$ to convergence in an inner loop before evaluating the conditional log-likelihood $\mathrm{LL}(\theta) = \sum_i \log \pi(a_i|s_i\,;\,\theta)$. The inner loop dominates computation.

In the tabular DDC reinterpretation of IQ-Learn, the Bellman equation is imposed as a soft constraint through the penalty term. The parameter is the $Q$-function itself, not the structural vector $\theta$. Because the transition matrix is known, the smoothed value function $V(s) = \sigma\log\left(\sum_a \exp\{Q(s,a)/\sigma\}\right)$ and the expected continuation value $\beta \sum_{s'} V(s')\,p(s'|s,a)$ are both explicit closed-form functions of $Q$. Each evaluation of \eqref{eq:penalizedMLE} requires one forward pass through these expressions. No fixed-point iteration appears. In continuous state-action spaces where $p$ is unknown, the original IQ-Learn paper replaces these exact computations with sampled next states and a separate policy network, as discussed in Section~\ref{sec:policy}.

The regularization parameter $\alpha$ governs how tightly the Bellman equation is enforced. As $\alpha \to 0$, the penalty dominates and forces the implied reward $\hat{r}_i$ toward zero at the expert-visited state-action pairs. Because $V$ depends on the entire $Q$-table through the logsumexp, this constraint propagates beyond the expert support, but the direct enforcement is only at states and actions observed in the data. As $\alpha \to \infty$, the penalty vanishes and the estimator reduces to an unconstrained conditional logit with no structural content.

\subsection{Comparison with CCP estimation}

CCP estimation avoids the Bellman fixed point by estimating choice probabilities nonparametrically in a first stage, then backing out structural parameters $\theta$ in a second stage. IQ-Learn avoids the fixed point by estimating $Q$ directly, then recovering the reward via $r(s,a) = Q(s,a) - \beta \sum_{s'} V(s')\, p(s'|s,a)$. Both approaches replace the nested fixed-point computation with a direct estimation step.

\subsection{Comparison with the ERM-IRL estimator of Kang et al.}

The ERM-IRL estimator of Kang et al.\ (2025), referred to as GLADIUS in the Rust and Rawat survey, shares the same high-level structure as the tabular IQ-Learn rewrite. Both objectives combine a conditional log-likelihood term with a Bellman-consistency penalty evaluated over expert observations.

The key difference is how the continuation value $\mathbb{E}\{V(s')|s,a\}$ is handled. Kang et al.\ work in the model-free setting where $p(s'|s,a)$ is unknown. Replacing the conditional expectation with a single observed draw $V(s'_i)$ introduces a double-sampling bias: the mean squared TD error $\rho_{TD}(Q) = \rho_{BE}(Q) + \beta^2\,\mathbb{E}\{[V(s') - \mathbb{E}\{V(s')|s,a\}]^2\}$ exceeds the mean squared Bellman error $\rho_{BE}(Q)$ by a variance term. Kang et al.\ correct this bias by introducing a dual variable $\zeta(s,a)$ that approximates the conditional expectation, turning the problem into a min-max optimization with provable convergence guarantees. In our tabular DDC rewrite, this issue does not arise because $\sum_{s'} V(s')\,p(s'|s,a)$ is computed exactly from the known transition matrix.

Kang et al.\ also note that occupancy-matching methods such as IQ-Learn may produce $Q$-functions that do not satisfy the Bellman equation, so that computing $r$ from $Q$ via the inverse Bellman operator may not yield a valid reward function. Their ERM-IRL estimator is designed to avoid this failure mode by directly enforcing Bellman consistency in the objective.

\subsection{Policy improvement in continuous spaces}
\label{sec:policy}

In the finite-state, finite-action setting, the optimal policy given $Q$ is $\pi(a|s) = \mathrm{softmax}(Q(s,a)/\sigma)$, computed exactly from \eqref{eq:ccp}. No separate policy update is needed.

Algorithm 1 in Garg et al.\ (2021) includes a policy improvement step because the paper targets continuous state-action spaces. In such spaces one cannot enumerate all actions to compute the logsumexp in \eqref{eq:smoothV}. A separate parameterized policy network is needed to approximate the softmax, and this network requires its own gradient updates. This is an implementation detail for the continuous case, not a conceptual element of the estimator.

\subsection{Identification}

Multiple reward functions can rationalize the same observed behavior. Given any function $h\!: S \to \mathbb{R}$, the ``shaped'' reward
\[
r'(s,a) = r(s,a) + h(s) - \beta \sum_{s'} h(s')\,p(s'|s,a)
\]
yields the same optimal policy as $r$ (Ng et al., 1999). This is what the IRL literature calls potential-based reward shaping, where $h$ is the potential function. Without further restrictions on rewards, only these equivalence classes are identified.

IQ-Learn does not resolve this fundamental non-identification. Instead, the chi-squared regularizer selects a particular element from the equivalence class. Expression \eqref{eq:penalizedMLE} penalizes the $\ell^2$ norm of the implied reward at expert-visited states. The minimizer of \eqref{eq:iqlearn} trades off conditional log-likelihood against reward magnitude, preferring $Q$-functions whose implied reward is small at expert-visited states. The numerical illustration in Section~\ref{sec:numerical} confirms this trade-off: IQ-Learn achieves a lower conditional log-likelihood than NFXP but recovers a smaller implied reward.

Different $f$-divergences induce different implicit penalties on the implied reward. The chi-squared divergence yields the quadratic penalty in \eqref{eq:penalizedMLE}. The Jensen-Shannon divergence yields a non-quadratic penalty that does not correspond to a standard norm. This explains the remark in Garg et al.\ (2021) that alternative divergences ``are not as readily interpretable'' as chi-squared. The choice of divergence determines which member of the equivalence class is selected. It does not affect the fit to observed actions.

\subsection{Summary}

In the finite discrete offline setting with known transitions, the IQ-Learn chi-squared objective can be rewritten as an expert-occupancy-weighted penalized conditional log-likelihood for the softmax policy induced by $Q$, with an $\ell^2$ penalty on the Bellman-implied reward. No inner Bellman solve is needed in this tabular setting because the smoothed value function is an explicit function of $Q$ and the continuation value is computed exactly from the transition matrix. The penalty selects from the equivalence class of observationally equivalent rewards by favoring solutions with small implied reward at expert-visited states.

\bigskip
\noindent\textbf{References}

\smallskip\noindent
Garg, D., Chakraborty, S., Cundy, C., Song, J., and Ermon, S.\ (2021). IQ-Learn: Inverse soft-Q learning for imitation. \textit{Advances in Neural Information Processing Systems 35}.

\smallskip\noindent
Kang, E.H., Yoganarasimhan, H., and Jain, L.\ (2025). An empirical risk minimization approach for offline inverse RL and dynamic discrete choice model. arXiv:2502.14131.

\smallskip\noindent
Ng, A., Harada, D., and Russell, S.\ (1999). Policy invariance under reward transformations: Theory and application to reward shaping. \textit{Proceedings of the 16th International Conference on Machine Learning}.

\end{document}
